\section{Introduction}
Self-evolving agents increasingly learn through external memory. A successful episode can become a reusable workflow. A failed episode can become a corrective reflection. Both objects may later be retrieved for a new task. Systems such as ReasoningBank make this loop explicit by converting trajectories into reasoning memories that influence future behavior \citep{ouyang2026reasoningbank}. Related systems store workflows or verbal feedback as persistent state \citep{wang2024awm,shinn2023reflexion}.

The usual abstraction treats memory content as the object that matters. This hides a second variable: how the memory was produced. A memory induced after success carries evidence that its reasoning process reached a valid outcome. A memory induced after failure carries evidence that the original process broke somewhere. ReasoningBank reflects this distinction in its implementation. Successful trajectories are summarized under a success-analysis objective. Failed trajectories are processed under a failure-analysis objective. The resulting memories can contain similar procedural advice while encoding different causal interpretations of the same task family.

A recent audit of self-evolving financial agents provides a strong empirical signal for this distinction \citep{li2026audit}. In its ReasoningBank analysis, success-derived retrievals have reported accuracy 0.931 over 101 cases. Failure-derived retrievals have reported accuracy 0.647 over 34 cases. The audit further reports ten persistent W$\rightarrow$W failures, where a task fails before evolution and continues to fail afterward. Every one of these persistent failures appears in the failure-derived retrieval stratum. Figure~\ref{fig:source} summarizes the released counts.

This association points to a causal question with direct implications for memory governance: does provenance itself change future behavior after semantic guidance is controlled? Task difficulty provides the strongest competing explanation. Hard tasks are more likely to fail during acquisition. Those same tasks can later retrieve memories derived from failure. A descriptive provenance gap can therefore emerge even when provenance has no causal effect. We resolve this identification problem with a matched provenance intervention.

The intervention begins with a future query and a retrieved failure-derived memory. We construct a success-derived counterpart that matches the guidance at the semantic level. The future query remains fixed. The agent scaffold remains fixed. Replacing only the memory provenance creates the causal contrast. A reverse swap tests the opposite direction. The resulting estimand measures whether failure-derived memory sustains error beyond the information carried by the memory itself.

The paper makes three contributions through separate evidence layers. First, we reanalyze the released financial-agent audit and show that the strongest provenance-conditioned signal is persistent failure. Second, we formalize provenance as an upstream state variable in memory construction. Third, we specify a matched-memory causal design that isolates the effect of provenance from task difficulty. These results motivate a provenance-aware representation for self-evolving memory systems. The source of a memory becomes part of the state that governs whether the agent should reuse it.
